\section{Funktory pochodne, Ext i Tor}\label{cat:sec:derived-functors}
Pracujemy w lokalnie małych kategoriach abelowych. Wykorzystujemy konstrukcje z rozdziałów~\ref{cat:sec:abelian-categories} i~\ref{cat:sec:projective-injective}. Wersję modułową teorii homologii i kohomologii omówiono w \atlasref{TOP-03::def:homologia-algebraiczna}{TOP-03: homologia} oraz \atlasref{TOP-03::def:kohomologia-algebraiczna}{TOP-03: kohomologia}.

\subsection{Homologia w kategorii abelowej}
\begin{definicja}[Cykle, brzegi i homologia]\label{cat:def:abelian-homology}
Niech $(C_\bullet,d)$ będzie kompleksem łańcuchowym, tj. rodziną obiektów $C_n$ i morfizmów $d_n:C_n\to C_{n-1}$ z $d_nd_{n+1}=0$. Dopuszczamy indeksy całkowite; dla kompleksu nieujemnie stopniowanego przyjmujemy $C_n=0$ dla $n<0$. Niech $z_n:Z_n\to C_n$ będzie jądrem $d_n$, a $b_n:B_n\to C_n$ obrazem $d_{n+1}$. Rozkład $d_{n+1}=b_ne_{n+1}$ z $e_{n+1}$ epi daje $d_nb_n=0$. Istnieje więc jedyny monomorfizm $j_n:B_n\to Z_n$ z $z_nj_n=b_n$; jest mono, ponieważ $b_n$ jest mono. Definiujemy
\[
 H_n(C_\bullet)=\operatorname{coker}(j_n)
\]
w sensie przeciwdziedziny kojądra $\pi_n:Z_n\to H_n(C_\bullet)$. Obiekty $Z_n,B_n$ nazywamy odpowiednio obiektami cykli i brzegów. Dla kompleksu kołańcuchowego stosujemy ten sam schemat z indeksami odwróconymi: $Z^n=\ker d^n$, $B^n=\operatorname{im}d^{n-1}$, $H^n=\operatorname{coker}(B^n\to Z^n)$, gdzie symbole oznaczają odpowiednie obiekty.
\end{definicja}
\begin{stwierdzenie}\label{cat:prop:homology-zero}
Kompleks jest dokładny w stopniu $n$ wtedy i tylko wtedy, gdy $H_n=0$ jako obiekt.
\end{stwierdzenie}
\begin{proof}
Kojądro $j_n$ jest zerowe dokładnie wtedy, gdy $j_n$ jest epi (stwierdzenie~\ref{cat:prop:mono-epi-zero-kernel}). Ponieważ $j_n$ jest mono, twierdzenie~\ref{cat:thm:abelian-factorization} daje równoważność z tym, że jest izomorfizmem, czyli że obraz i jądro określają ten sam podobiekt $C_n$. To jest dokładność.
\end{proof}
\begin{twierdzenie}[Funktorialność i niezmienniczość homotopijna]\label{cat:thm:abelian-homology-functor}
Dla ustalonego wyboru jąder, obrazów i kojąder $H_n$ jest funktorem addytywnym z kategorii kompleksów łańcuchowych do kategorii abelowej. Morfizmy homotopijne indukują ten sam morfizm homologii. Analogiczne twierdzenie obowiązuje dla kohomologii.
\end{twierdzenie}
\begin{proof}
Dla morfizmu kompleksów $f:C_\bullet\to D_\bullet$ mamy $d_n^D f_nz_n^C=f_{n-1}d_n^Cz_n^C=0$. Istnieje jedyny $f_Z:Z_n^C\to Z_n^D$ z $z_n^Df_Z=f_nz_n^C$. Ponadto
\[
 f_nb_n^Ce_{n+1}^C=f_nd_{n+1}^C=d_{n+1}^Df_{n+1}
 =b_n^De_{n+1}^Df_{n+1}.
\]
Jeśli $q$ jest kojądrem $b_n^D$, to $q f_nb_n^Ce_{n+1}^C=0$, a $e_{n+1}^C$ jest epi, więc $qf_nb_n^C=0$. Ponieważ $b_n^D=\ker q$ (twierdzenie~\ref{cat:thm:abelian-factorization}), istnieje jedyny $f_B:B_n^C\to B_n^D$ z $b_n^Df_B=f_nb_n^C$. Z monomorficzności $z_n^D$ wynika $f_Zj_n^C=j_n^Df_B$. Zatem $\pi_n^Df_Zj_n^C=0$, co daje jedyny morfizm $H_n(f)$ z $H_n(f)\pi_n^C=\pi_n^Df_Z$. Jednoznaczność konstrukcji na jądrach i kojądrach daje zachowywanie identyczności i złożeń. Addytywność wynika z równości $(f+g)_Z=f_Z+g_Z$ (składamy z mono $z_n^D$) i $(H_n(f)+H_n(g))\pi_n^C=\pi_n^D(f_Z+g_Z)$; skracamy epi $\pi_n^C$.

Jeżeli $f-g=d^Dh+hd^C$, to $(f_n-g_n)z_n^C=d_{n+1}^Dh_nz_n^C=b_n^De_{n+1}^Dh_nz_n^C$. Po skróceniu mono $z_n^D$ otrzymujemy $f_Z-g_Z=j_n^De_{n+1}^Dh_nz_n^C$. Składając z $\pi_n^D$, otrzymujemy zero, więc $(H_n(f)-H_n(g))\pi_n^C=0$. Epimorficzność $\pi_n^C$ kończy dowód. Dla kompleksów kołańcuchowych zamieniamy $C_n$ na $C^{-n}$ i stosujemy dowód bez zmian.
\end{proof}

\subsection{Lewe i prawe funktory pochodne}
\begin{definicja}[Funktory pochodne]\label{cat:def:derived-functors}
Niech $F:\mathcal A\to\mathcal B$ będzie funktorem addytywnym między kategoriami abelowymi. Jeżeli $\mathcal A$ ma dostatecznie wiele obiektów rzutowych, wybieramy rezolucję rzutową $P_\bullet(X)\to X$ i kładziemy
\[
 L_nF(X)=H_n(F(P_\bullet(X)))\quad(n\ge0).
\]
Jeżeli $\mathcal A$ ma dostatecznie wiele obiektów injektywnych, wybieramy rezolucję injektywną $X\to I^\bullet(X)$ i kładziemy
\[
 R^nF(X)=H^n(F(I^\bullet(X)))\quad(n\ge0).
\]
W obu wzorach pomijamy augmentację lub koaugmentację przed zastosowaniem homologii. Addytywność $F$ zapewnia zachowanie morfizmów zerowych, zatem obrazy kompleksów pozostają kompleksami. Przy wyborach dotyczących wszystkich obiektów pracujemy w ustalonym uniwersum, w którym można dokonać takich wyborów.
\end{definicja}
\begin{twierdzenie}[Dobrze określone funktory pochodne]\label{cat:thm:derived-well-defined}
Powyższe konstrukcje wyznaczają funktory addytywne, niezależne od wyboru rezolucji z dokładnością do kanonicznych izomorfizmów naturalnych zgodnych z kolejną zmianą wyborów.
\end{twierdzenie}
\begin{proof}
Dla $u:X\to Y$ wybieramy podniesienie $f:P_\bullet(X)\to P_\bullet(Y)$ z twierdzenia~\ref{cat:thm:resolution-comparison}. Definiujemy $L_nF(u)=H_n(Ff)$. Dwa podniesienia są homotopijne; zastosowanie addytywnego $F$ do równań homotopii zachowuje wszystkie sumy i złożenia, więc ich obrazy pozostają homotopijne. Twierdzenie~\ref{cat:thm:abelian-homology-functor} daje tę samą strzałkę homologii. Złożenie podniesień podnosi złożenie morfizmów, identyczność podnosi identyczność, a suma podniesień podnosi sumę; stąd funktorialność i addytywność.
Dla drugiego wyboru rezolucji $Q_\bullet(X)$ podniesienie $\id_X$ daje morfizm $P_\bullet(X)\to Q_\bullet(X)$. Podniesienie w drugą stronę jest jego odwrotnością z dokładnością do homotopii na mocy wniosku~\ref{cat:cor:resolution-homotopy-equivalence}. Po zastosowaniu $F$ i $H_n$ otrzymujemy izomorfizm $\alpha_X$. Dla $u:X\to Y$ oba złożenia w kwadracie naturalności przed zastosowaniem $F$ i $H_n$ podnoszą $u$ z $P_\bullet(X)$ do $Q_\bullet(Y)$, więc są homotopijne. Stąd kwadrat jest przemienny. Przy trzecim wyborze rezolucji złożenie dwóch izomorfizmów odpowiada kolejnemu podniesieniu $\id_X$, więc jest równe bezpośredniemu izomorfizmowi. Dla $R^nF$ stosujemy wniosek~\ref{cat:cor:injective-comparison} i kohomologiczną wersję twierdzenia~\ref{cat:thm:abelian-homology-functor}; wszystkie argumenty o podniesieniach zastępujemy przedłużeniami.
\end{proof}
\begin{stwierdzenie}[Stopień zerowy i zanikanie]\label{cat:prop:derived-degree-zero}
Jeżeli $F$ jest prawostronnie dokładny, to $L_0F\cong F$ naturalnie. Jeżeli jest lewostronnie dokładny, to $R^0F\cong F$ naturalnie. Dla obiektu rzutowego $P$ mamy $L_nF(P)=0$ dla $n>0$, a dla injektywnego $I$ mamy $R^nF(I)=0$ dla $n>0$. Jeśli $\operatorname{pd}X\le r$, to $L_nF(X)=0$ dla $n>r$; dualnie dla wymiaru injektywnego.
\end{stwierdzenie}
\begin{proof}
W rezolucji rzutowej $P_1\xrightarrow{d_1}P_0\xrightarrow{\epsilon}X\to0$ augmentacja jest kojądrem $d_1$: ich obraz i jądro są równe, a $\epsilon$ jest epi, więc własność uniwersalna kojądra $\ker\epsilon$ jest także własnością kojądra $d_1$. Prawostronna dokładność oznacza zachowywanie kojąder (twierdzenie~\ref{cat:thm:exact-functor-kernels}); zatem $F\epsilon$ jest kojądrem $Fd_1$. Ponieważ w kompleksie bez augmentacji $Z_0=FP_0$, mamy $L_0F(X)\cong FX$. Dla podniesienia $f$ morfizmu $u$ zachodzi $F\epsilon_Y Ff_0=Fu F\epsilon_X$; własność kojądra daje naturalność. Dualnie $F\eta$ jest jądrem $Fd^0$ dla lewostronnie dokładnego $F$, co daje naturalne $R^0F\cong F$.
Dla $P$ wybieramy rezolucję z identycznością w stopniu zerowym i zerami w dodatnich stopniach. Funktor addytywny zachowuje obiekty zerowe, więc wyższe homologie są zerowe. Dla skończonej rezolucji długości $r$ obiekty w stopniach $n>r$ są zerowe i otrzymujemy to samo zanikanie. Argument injektywny jest dualny.
\end{proof}
\begin{stwierdzenie}\label{cat:prop:exact-derived-zero}
Jeżeli $F$ jest dokładny, wszystkie jego istniejące funktory pochodne dodatniego stopnia znikają.
\end{stwierdzenie}
\begin{proof}
Funktor dokładny zachowuje dokładność ciągów, w szczególności rezolucji z augmentacją lub koaugmentacją. Po usunięciu tych strzałek dokładność w dodatnich stopniach pozostaje bez zmian. Stwierdzenie~\ref{cat:prop:homology-zero} daje tezę.
\end{proof}

\subsection{Długi ciąg dokładny homologii}
\begin{twierdzenie}[Długi ciąg dokładny]\label{cat:thm:abelian-homology-les}
Niech $0\to A_\bullet\xrightarrow{i}B_\bullet\xrightarrow{p}C_\bullet\to0$ będzie ciągiem kompleksów łańcuchowych dokładnym w każdym stopniu. Istnieją naturalne morfizmy łączące $\partial_n:H_n(C_\bullet)\to H_{n-1}(A_\bullet)$, dla których ciąg
\[
\cdots\to H_n(A_\bullet)\xrightarrow{H_n(i)}H_n(B_\bullet)
\xrightarrow{H_n(p)}H_n(C_\bullet)
\xrightarrow{\partial_n}H_{n-1}(A_\bullet)\to\cdots
\]
jest dokładny. Dla kompleksów kołańcuchowych morfizm łączący zwiększa stopień o jeden.
\end{twierdzenie}
\begin{proof}
Piszemy $z_n^X:Z_n^X\to X_n$ dla jądra różniczki, $\pi_n^X:Z_n^X\to H_n(X)$ dla projekcji na homologię, a $d_{n+1}^X=z_n^X j_n^X e_{n+1}^X$ dla rozkładu różniczki przez cykle i brzegi, gdzie $e_{n+1}^X$ jest epi.

\emph{Konstrukcja.}
Tworzymy iloczyn włóknisty $D=B_n\times_{C_n} Z_n^C$ z projekcjami $t:D\to B_n$, $r:D\to Z_n^C$. Morfizm $r$ jest epi. Ponieważ $p_{n-1}d_n^Bt=d_n^Cz_n^Cr=0$, istnieje jedyny $a:D\to A_{n-1}$ z $i_{n-1}a=d_n^Bt$. Mamy $i_{n-2}d_{n-1}^Aa=d_{n-1}^Bd_n^Bt=0$, więc $d_{n-1}^Aa=0$. Istnieje jedyny $v:D\to Z_{n-1}^A$ z $z_{n-1}^Av=a$. Jeśli $rh=0$, to $p_nth=0$, więc $th=i_nu$ dla pewnego $u$ i $z_{n-1}^Avh=d_n^Au$. Stąd $vh=j_{n-1}^Ae_n^Au$ i $\pi_{n-1}^Avh=0$. W szczególności $\pi_{n-1}^Av$ zeruje się na jądrze $r$ i rozkłada się jednoznacznie przez epi $r$ jako $\lambda r$, gdzie $\lambda:Z_n^C\to H_{n-1}(A)$.

Aby $\lambda$ rozkładało się przez $\pi_n^C$, wystarczy wykazać $\lambda j_n^C=0$. Tworzymy morfizm $w:B_{n+1}\to D$ o składowych $d_{n+1}^B$ i $j_n^C e_{n+1}^C p_{n+1}$; warunek iloczynu włóknistego wynika z przemienności $p$ z różniczkami. Wówczas $vw=0$, ponieważ $i_{n-1}z_{n-1}^Avw=d_n^Bd_{n+1}^B=0$. Zatem $\lambda j_n^C e_{n+1}^C p_{n+1}=0$. Ostatnie dwa morfizmy są epi, więc $\lambda j_n^C=0$. Otrzymujemy jedyny morfizm $\partial_n$ z
\begin{equation}\label{cat:eq:homology-connecting}
 \partial_n\pi_n^Cr=\pi_{n-1}^Av.
\end{equation}

\emph{Podnoszenie lokalne.}
W dowodzie dokładności stosujemy lemat~\ref{cat:lem:abelian-local-lift}: gdy chcemy podnieść morfizm przez epi, możemy zastąpić jego dziedzinę $T$ przez dziedzinę epi $T'\to T$. Wszystkie wcześniej uzyskane morfizmy składamy wówczas z tym epi. Piszemy równości bez oznaczania kolejnych zmian dziedziny; każda równość po zmianie ma wspólną nową dziedzinę. Kryterium~\ref{cat:lem:exactness-local-criterion} pozwala w ten sposób dowodzić dokładności. W szczególności każdy morfizm do $H_n(X)$ można lokalnie podnieść do $Z_n^X$, a cykl o zerowej klasie homologii można lokalnie zapisać jako różniczkę: jest to zastosowanie lematu do dokładnego ciągu $X_{n+1}\to Z_n^X\to H_n(X)\to0$, którego dokładność wynika z definicji obrazu i kojądra.

\emph{Zerowanie złożeń i dokładność.}
Przez cykl w poniższych rachunkach rozumiemy morfizm $x:T\to X_n$ z $d_n^Xx=0$, a przez $[x]$ złożenie jego jednoznacznego rozkładu przez $z_n^X$ z $\pi_n^X$. Każdy morfizm homologii jest lokalnie takiej postaci. Mamy $H_n(p)H_n(i)=0$ z $pi=0$. Dla cyklu $b$ w $B_n$ równanie~\eqref{cat:eq:homology-connecting} daje $\partial_n[p_nb]=0$, gdyż można użyć podniesienia $b$ o zerowej różniczce. Dla lokalnego podniesienia $b$ cyklu $c$ z $C_n$ mamy $p_nb=c$, $d_n^Bb=i_{n-1}a$ i $\partial_n[c]=[a]$. Zatem $H_{n-1}(i)\partial_n[c]=[d_n^Bb]=0$. Te złożenia zerują się na wszystkich morfizmach, gdyż epi użyte do zmian dziedziny można skrócić.

Jeśli $b$ jest cyklem w $B_n$ i $H_n(p)[b]=0$, to lokalnie $p_nb=d_{n+1}^Cc'$ dla $c':T'\to C_{n+1}$. Podnosimy $c'$ przez epi $p_{n+1}$ do $b'$. Wtedy $p_n(b-d_{n+1}^Bb')=0$, więc $b-d_{n+1}^Bb'=i_na$. Monomorficzność $i_{n-1}$ daje $d_n^Aa=0$. Stąd $[b]=H_n(i)[a]$.

Jeśli $c$ jest cyklem w $C_n$ i $\partial_n[c]=0$, podnosimy $c$ lokalnie do $b$ przez $p_n$. Mamy $d_n^Bb=i_{n-1}a$ z $a$ cyklem w $A_{n-1}$ i $[a]=0$. Lokalnie $a=d_n^Aa'$; wtedy $b-i_na'$ jest cyklem w $B_n$ i $p_n(b-i_na')=c$. Zatem $[c]=H_n(p)[b-i_na']$.

Jeśli $a$ jest cyklem w $A_{n-1}$ i $H_{n-1}(i)[a]=0$, to lokalnie $i_{n-1}a=d_n^Bb$ dla pewnego $b$. Wtedy $c=p_nb$ jest cyklem, gdyż $d_n^Cc=p_{n-1}i_{n-1}a=0$. Równanie~\eqref{cat:eq:homology-connecting} daje $\partial_n[c]=[a]$. Te trzy rachunki wraz z kryterium lokalnego podnoszenia dowodzą dokładności we wszystkich stopniach.

\emph{Naturalność.}
Dla morfizmu krótkich ciągów kompleksów $(f_A,f_B,f_C)$ lokalne podniesienie $b$ cyklu $c$ przechodzi na $f_{B,n}b$, podniesienie $f_{C,n}c$. Jeśli $d_n^Bb=i_{n-1}a$, to $d_n^{B'}f_{B,n}b=i'_{n-1}f_{A,n-1}a$. Z definicji morfizmu łączącego wynika zatem
$\partial'_nH_n(f_C)[c]=H_{n-1}(f_A)\partial_n[c]$.
Każda klasa ma lokalnego reprezentanta, więc po skróceniu epi otrzymujemy równość morfizmów. Pozostałe kwadraty są przemienne przez funktorialność homologii. Odwrócenie indeksów daje wersję kołańcuchową.
\end{proof}

\begin{twierdzenie}[Długie ciągi funktorów pochodnych]\label{cat:thm:derived-les}
Dla addytywnego funktora $F:\mathcal A\to\mathcal B$ oraz krótkiego ciągu dokładnego $0\to X\to Y\to Z\to0$ istnieją naturalne długie ciągi
\[
 \cdots\to L_nF(X)\to L_nF(Y)\to L_nF(Z)
 \to L_{n-1}F(X)\to\cdots\to L_0F(Z)\to0
\]
przy dostatecznej liczbie obiektów rzutowych, a przy dostatecznej liczbie injektywnych
\[
 0\to R^0F(X)\to R^0F(Y)\to R^0F(Z)
 \to R^1F(X)\to R^1F(Y)\to\cdots .
\]
\end{twierdzenie}
\begin{proof}
Lemat o podkowie (twierdzenie~\ref{cat:thm:horseshoe}) daje krótki ciąg rezolucji rozkładalny w każdym stopniu. Addytywny $F$ zachowuje biprodukty, więc po jego zastosowaniu ciąg pozostaje dokładny w każdym stopniu, niezależnie od dokładności $F$. Twierdzenie~\ref{cat:thm:abelian-homology-les} daje długi ciąg, a twierdzenie~\ref{cat:thm:derived-well-defined} identyfikuje jego obiekty i strzałki niełączące z ustalonymi funktorami pochodnymi.

Dla naturalności potrzebujemy zgodnych morfizmów między ciągami rezolucji z lematu o podkowie. Oznaczmy w stopniu $n$ składniki przez $P_n,T_n=P_n\oplus Q_n,Q_n$ i analogicznie z primem. W stopniu $0$ wybieramy podniesienia $f_P,f_Q$ morfizmów obiektów skrajnych przez augmentacje. Na środkowym obiekcie szukamy macierzy
\[
 f_T=\begin{pmatrix}f_P&h\\0&f_Q\end{pmatrix}.
\]
Jeśli na razie przyjmiemy $h=0$, różnica między wymaganym złożeniem z augmentacją a aktualnym złożeniem znika na $P_0$ i po projekcji do obiektu skrajnego prawego. Na $Q_0$ rozkłada się więc przez obiekt skrajny lewy docelowego ciągu. Rzutowość $Q_0$ pozwala podnieść tę różnicę przez augmentację $P'_0$; otrzymane $h:Q_0\to P'_0$ usuwa różnicę. Tak powstały przemienny diagram indukuje morfizm krótkich ciągów jąder. W następnym stopniu stosujemy ten sam argument do tych jąder i epimorfizmów użytych w konstrukcji podkowy. Indukcja daje morfizm krótkich ciągów rezolucji zgodny z augmentacjami. Twierdzenie o naturalności homologii daje naturalność strzałek łączących. Stosując konstrukcję do identyczności krótkiego ciągu i dwóch wyborów podków, otrzymujemy zgodność strzałek łączących z kanonicznymi izomorfizmami zmiany rezolucji. Wersję injektywną otrzymujemy odwracając strzałki w kategorii przeciwnej i stosując wersję kołańcuchową homologii.
\end{proof}

\subsection{Funktory Ext i Tor}
\begin{definicja}\label{cat:def:ext-tor}
Niech $R$ będzie pierścieniem z jedynką. Dla lewych $R$-modułów $M,N$ definiujemy
\[
 \operatorname{Ext}_R^n(M,N)=H^n(\Hom_R(P_\bullet(M),N)).
\]
Jest to prawy funktor pochodny kontrawariantnego funktora $\Hom_R(-,N)$ rozpatrywanego na kategorii przeciwnej; rezolucja rzutowa w kategorii modułów staje się tam rezolucją injektywną. Dla prawego $R$-modułu $U$ i lewego $R$-modułu $M$ definiujemy
\[
 \operatorname{Tor}_n^R(U,M)=H_n(U\otimes_R P_\bullet(M)).
\]
Jest to lewy funktor pochodny funktora $U\otimes_R-$ o wartościach w grupach abelowych. Nie zakładamy przemienności $R$. Obie konstrukcje są niezależne od wyboru rezolucji na mocy twierdzenia~\ref{cat:thm:derived-well-defined}.
\end{definicja}
\begin{stwierdzenie}\label{cat:prop:ext-tor-basic}
Mamy naturalne izomorfizmy $\operatorname{Ext}_R^0(M,N)\cong\Hom_R(M,N)$ oraz $\operatorname{Tor}_0^R(U,M)\cong U\otimes_R M$. Ext jest kontrawariantny w pierwszej i kowariantny w drugiej zmiennej, a Tor kowariantny w obu. Przy rzutowym $M$ oba funktory znikają w dodatnich stopniach. Dla injektywnego $N$ również $\operatorname{Ext}_R^n(M,N)=0$ dla $n>0$.
\end{stwierdzenie}
\begin{proof}
Funktor $\Hom_R(-,N)$ jest lewostronnie dokładny na kategorii przeciwnej, natomiast $U\otimes_R-$ jest prawostronnie dokładny, gdyż jest lewym sprzężonym funktora $\Hom_{\mathbb Z}(U,-)$ z lewym działaniem $(rf)(u)=f(ur)$. Własność sprzężenia wynika z bijekcji między homomorfizmami $U\otimes_R M\to D$ i $M\to\Hom_{\mathbb Z}(U,D)$, danej przez $g\mapsto(m\mapsto(u\mapsto g(u\otimes m)))$; relacja balansowania $ur\otimes m=u\otimes rm$ odpowiada $R$-liniowości. Twierdzenie~\ref{cat:thm:adjoint-exactness} i stwierdzenie~\ref{cat:prop:derived-degree-zero} dają izomorfizmy stopnia zerowego oraz zanikanie przy rzutowym $M$. Morfizmy w pozostałej zmiennej działają stopniowo przez postkompozycję w Hom lub przez iloczyn tensorowy z identycznością; są to morfizmy kompleksów. Komutują z morfizmami pochodzącymi z rezolucji, gdyż składanie homomorfizmów jest łączne, a iloczyn tensorowy jest bifunktorem. Stąd otrzymujemy bifunktory. Przy injektywnym $N$ funktor $\Hom_R(-,N)$ jest dokładny na mocy twierdzenia~\ref{cat:thm:projective-characterization}, więc wyższe funktory znikają.
\end{proof}
\begin{przyklad}[Grupy cykliczne]\label{cat:ex:ext-tor-cyclic}
Dla grupy abelowej $N$ i $m\ge2$ stosujemy rezolucję $0\to\mathbb Z\xrightarrow{m}\mathbb Z\to\mathbb Z/m\mathbb Z\to0$ z przykładu~\ref{cat:ex:cyclic-projective-resolution}. Utożsamienia $\Hom_{\mathbb Z}(\mathbb Z,N)\cong N$, $f\mapsto f(1)$, i $N\otimes_{\mathbb Z}\mathbb Z\cong N$, $x\otimes k\mapsto kx$, przekształcają odpowiednie różniczki na mnożenie przez $m$. Stąd
\[
\begin{aligned}
 \operatorname{Ext}_{\mathbb Z}^1(\mathbb Z/m\mathbb Z,N)&\cong N/mN,\\
 \operatorname{Tor}_1^{\mathbb Z}(N,\mathbb Z/m\mathbb Z)&\cong\{x\in N:mx=0\}.
\end{aligned}
\]
Wyższe stopnie $n\ge2$ znikają. Jeśli $N=\mathbb Z/k\mathbb Z$ i $g=\gcd(m,k)$, to jądro mnożenia przez $m$ jest generowane przez klasę $k/g$ o rzędzie $g$, a obraz przez klasę $g$: z tożsamości Bézouta wynika $g\in m\mathbb Z+k\mathbb Z$. Iloraz przez obraz ma więc rząd $g$ i jest cykliczny. Obie grupy są izomorficzne z $\mathbb Z/g\mathbb Z$, co nie oznacza, że Ext i Tor są tym samym funktorem.
\end{przyklad}

\subsection{Kompleksy podwójne i zrównoważoność}
\begin{lemat}[Porównanie z kompleksem totalnym]\label{cat:lem:total-acyclic}
Niech $D_{p,q}$, $p,q\ge0$, będzie podwójnym kompleksem grup abelowych z różniczkami poziomą $d_h$ i pionową $d_v$, które komutują i osobno mają kwadrat zero. Kompleks totalny ma $\operatorname{Tot}_n D=\bigoplus_{p+q=n}D_{p,q}$ i różniczkę $d=d_h+(-1)^p d_v$ na $D_{p,q}$.
Jeżeli wszystkie poziome wiersze mają dokładne augmentacje $D_{0,q}\to E_q\to0$, zgodne z różniczkami pionowymi, to naturalny morfizm $\operatorname{Tot}D\to E_\bullet$ indukuje izomorfizmy homologii. Analogicznie dla kompleksów podwójnych kołańcuchowych $D^{p,q}$ z różniczkami zwiększającymi indeksy, jeśli pionowe kolumny mają dokładne koaugmentacje $0\to E^p\to D^{p,0}$, zgodne z różniczkami poziomymi, to morfizm $E^\bullet\to\operatorname{Tot}D$ indukuje izomorfizmy kohomologii.
\end{lemat}
\begin{proof}
Z przemienności różniczek wynika $d^2=0$: dwa mieszane składniki na $D_{p,q}$ mają współczynniki $(-1)^p$ oraz $(-1)^{p-1}$; w wersji kołańcuchowej odpowiednio $(-1)^p$ i $(-1)^{p+1}$. W obu przypadkach się znoszą. Wszystkie sumy w ustalonym stopniu są skończone.

W wersji łańcuchowej projekcja na $p=0$ i augmentacja dają surjektywny w każdym stopniu morfizm kompleksów. Jego jądro jest totalizacją podwójnego kompleksu $K$ z $K_{p,q}=D_{p,q}$ dla $p>0$ oraz $K_{0,q}=\ker(D_{0,q}\to E_q)$. Wiersze $K$ są dokładne. Pokażemy, że każdy cykl $x$ stopnia $n$ w $\operatorname{Tot}K$ jest brzegiem. Niech $q$ będzie największym indeksem pionowym niezerowej składowej $x_{p,q}$, gdzie $p+q=n$. Jeśli $p>0$, równanie $dx=0$ w składowej $(p-1,q)$ daje $d_hx_{p,q}=0$, gdyż nie ma składowej o indeksie pionowym $q+1$. Dla $p=0$ ten sam warunek zachodzi przez przyjęcie $K_{-1,q}=0$. Dokładność wiersza pozwala wybrać $y_{p+1,q}$ z $d_hy_{p+1,q}=x_{p,q}$. Odejmując $dy_{p+1,q}$, usuwamy tę składową i zmieniamy co najwyżej składową o indeksie pionowym $q-1$. Powtarzamy aż do $q=0$, otrzymując $x=d(\sum y)$. Jądro jest zatem acykliczne. Długi ciąg dokładny z twierdzenia~\ref{cat:thm:abelian-homology-les} daje tezę.

W wersji kołańcuchowej koaugmentacja daje monomorfizm kompleksów. Jego iloraz jest totalizacją podwójnego kompleksu $K$ o dokładnych kolumnach pionowych: w stopniu $q=0$ zastępujemy $D^{p,0}$ ilorazem przez obraz $E^p$, a w wyższych stopniach nic nie zmieniamy. Dla cocyklu $x$ stopnia $n$ wybieramy najmniejszy indeks poziomy $p$ niezerowej składowej $x^{p,q}$, $p+q=n$. Równanie $dx=0$ w składowej $(p,q+1)$ daje $d_vx^{p,q}=0$, ponieważ nie ma składowej o indeksie $p-1$. Jeśli $q=0$, dokładność kolumny wymusza $x^{p,0}=0$. Jeśli $q>0$, istnieje $y^{p,q-1}$ z $d_vy^{p,q-1}=(-1)^p x^{p,q}$. Odejmując $dy^{p,q-1}$, usuwamy składową o indeksie $p$, zmieniając co najwyżej następną o indeksie $p+1$. Po skończenie wielu krokach otrzymujemy zero. Iloraz jest acykliczny, więc długi ciąg dokładny kohomologii daje izomorfizm. Obie konstrukcje są naturalne względem morfizmów podwójnych kompleksów zgodnych z augmentacjami.
\end{proof}
\begin{uwaga}
W lemacie można zamienić role wierszy i kolumn. Przy transpozycji $(p,q)\mapsto(q,p)$ mnożenie składowej przez $(-1)^{pq}$ daje izomorfizm kompleksów totalnych: dla różniczki poziomej zmiana znaku wynosi $(-1)^q$, a dla pionowej $(-1)^p$. Dotyczy to zarówno zwiększania, jak i zmniejszania indeksów.
\end{uwaga}

\begin{twierdzenie}[Zrównoważoność Ext]\label{cat:thm:ext-balanced}
Dla lewych $R$-modułów $M,N$ i dowolnej rezolucji injektywnej $N\to I^\bullet$ mamy naturalny izomorfizm
\[
 \operatorname{Ext}_R^n(M,N)\cong H^n(\Hom_R(M,I^\bullet)).
\]
\end{twierdzenie}
\begin{proof}
Wybieramy rezolucję rzutową $P_\bullet\to M$. Rozważamy $D^{p,q}=\Hom_R(P_p,I^q)$ z różniczkami $d_h(f)=f d_{p+1}^P$ i $d_v(f)=d_I^q f$. Mają one kwadrat zero i komutują przez łączność składania. Dla ustalonego $p$ funktor $\Hom_R(P_p,-)$ jest dokładny przez rzutowość $P_p$. Kolumna ma więc dokładną koaugmentację $\Hom_R(P_p,N)\to D^{p,0}$. Dla ustalonego $q$ dokładność $\Hom_R(-,I^q)$ przez injektywność $I^q$ daje dokładną koaugmentację wiersza $\Hom_R(M,I^q)\to D^{0,q}$. Obie koaugmentacje są zgodne z pozostałymi różniczkami przez przemienność złożeń. Lemat~\ref{cat:lem:total-acyclic} oraz transpozycja dają izomorfizmy
\[
 H^n(\Hom_R(P_\bullet,N))
 \longrightarrow H^n(\operatorname{Tot}D)
 \longleftarrow H^n(\Hom_R(M,I^\bullet)).
\]
Złożenie pierwszego z odwrotnością drugiego daje żądany izomorfizm. Dla morfizmu w pierwszej zmiennej wybieramy podniesienie między rezolucjami rzutowymi i stosujemy prekompozycję; dla morfizmu w drugiej zmiennej wybieramy przedłużenie między rezolucjami injektywnymi i stosujemy postkompozycję. Otrzymujemy morfizmy podwójnych kompleksów zgodne z obiema koaugmentacjami. Przemienność kwadratów indukowanych w kohomologii daje naturalność. Dla identyczności obiektów i zmiany rezolucji ta sama przemienność daje zgodność z kanonicznymi izomorfizmami z twierdzenia~\ref{cat:thm:derived-well-defined}. Zatem otrzymany izomorfizm nie zależy od wyboru rezolucji.
\end{proof}

\begin{definicja}[Moduł płaski]\label{cat:def:flat}
Lewy moduł $F$ jest płaski, jeśli funktor $-\otimes_R F$ na prawych modułach jest dokładny; prawy moduł $G$ jest płaski, jeśli $G\otimes_R-$ jest dokładny na lewych modułach. Por. \atlasref{AA-18::aa:def:flat-module}{AA-18: moduły płaskie}.
\end{definicja}
\begin{stwierdzenie}\label{cat:prop:projective-flat}
Każdy moduł rzutowy jest płaski.
\end{stwierdzenie}
\begin{proof}
Dla modułów wolnych $R^{(S)}$ funktor iloczynu tensorowego z nimi utożsamia się z sumą prostą $S$ kopii drugiego modułu, która zachowuje dokładność: cykle i obrazy wyznacza się współrzędnie, a element ma skończenie wiele niezerowych współrzędnych. Utożsamienie wynika z własności uniwersalnej iloczynu tensorowego; por. \atlasref{AA-18::aa:tw:tensor-direct-sum}{AA-18: iloczyn tensorowy i sumy proste} oraz \atlasref{AA-18::aa:tw:tensor-unit}{AA-18: iloczyn tensorowy z pierścieniem}. Moduł rzutowy jest składnikiem sumy prostej modułu wolnego (dla prawych modułów stosujemy twierdzenie o lewych modułach nad $R^{\op}$). Iloczyn tensorowy rozkłada się zgodnie z tą sumą prostą. Dokładność sumy dwóch ciągów implikuje dokładność każdego składnika, ponieważ jądra i obrazy rozkładają się na sumy prostych odpowiednich jąder i obrazów. Stąd płaskość; por. również \atlasref{AA-18::aa:tw:projective-flat}{AA-18: moduły rzutowe są płaskie}.
\end{proof}
\begin{twierdzenie}[Zrównoważoność Tor i rezolucje płaskie]\label{cat:thm:tor-balanced}
Niech $U$ będzie prawym, a $M$ lewym $R$-modułem. Jeżeli $G_\bullet\to U$ jest dokładną rezolucją przez płaskie prawe moduły, a $F_\bullet\to M$ dokładną rezolucją przez płaskie lewe moduły, to naturalnie
\[
 \operatorname{Tor}_n^R(U,M)\cong H_n(G_\bullet\otimes_R M)
 \cong H_n(U\otimes_R F_\bullet).
\]
W szczególności dowolna rezolucja rzutowa prawego modułu $U$ może służyć do obliczania Tor.
\end{twierdzenie}
\begin{proof}
Najpierw rozważamy dwie rezolucje rzutowe $Q_\bullet\to U$, $P_\bullet\to M$. Kompleks podwójny $D_{p,q}=Q_p\otimes_R P_q$ ma komutujące różniczki $d_h=d^Q\otimes1$, $d_v=1\otimes d^P$; komutują one dzięki bifunktorialności iloczynu tensorowego. Dokładność wierszy augmentowanych do $U\otimes_R P_q$ wynika z płaskości $P_q$, a kolumn augmentowanych do $Q_p\otimes_R M$ z płaskości $Q_p$. Lemat~\ref{cat:lem:total-acyclic} i transpozycja dają izomorfizmy
\[
 H_n(U\otimes_R P_\bullet)
 \longleftarrow H_n(\operatorname{Tot}D)
 \longrightarrow H_n(Q_\bullet\otimes_R M).
\]
Podniesienia morfizmów $U,M$ do rezolucji rzutowych indukują morfizmy kompleksów podwójnych zgodne z augmentacjami; naturalność lematu daje naturalność izomorfizmu. Niezależność od wyboru rezolucji wynika z zastosowania tej samej naturalności do podniesień identyczności.

Następnie stosujemy argument do $D_{p,q}=Q_p\otimes_R F_q$. Dokładność augmentowanych wierszy wynika z płaskości $F_q$, a kolumn z płaskości $Q_p$. Otrzymujemy izomorfizm $H_n(U\otimes_R F_\bullet)\cong H_n(Q_\bullet\otimes_R M)$, a ten drugi obiekt oblicza Tor na mocy pierwszej części dowodu. Zamiana $Q,P$ na $G,P$ daje analogiczny wynik dla $H_n(G_\bullet\otimes_R M)$. Nie zakładamy istnienia podniesień dowolnych morfizmów między rezolucjami płaskimi: do porównania służą podniesienia z rezolucji rzutowych (twierdzenie~\ref{cat:thm:resolution-comparison}). Dla dowolnego takiego podniesienia otrzymujemy przemienny diagram z augmentacjami w kompleksach podwójnych po użyciu rezolucji rzutowej w drugiej zmiennej. Stąd izomorfizm jest kanoniczny i zgodny z morfizmami rezolucji płaskich, gdy takie morfizmy są dane. Naturalne działanie na Tor jest zawsze określone przez rezolucje rzutowe.
\end{proof}
\begin{wniosek}\label{cat:cor:tor-flat}
Jeżeli $M$ lub $U$ jest płaski, to $\operatorname{Tor}_n^R(U,M)=0$ dla $n>0$.
\end{wniosek}
\begin{proof}
Dla płaskiego modułu można wybrać rezolucję płaską długości zero z identycznością jako augmentacją. Twierdzenie~\ref{cat:thm:tor-balanced} pozwala obliczyć Tor przez tę rezolucję. W dodatnich stopniach kompleks jest zerowy, więc jego homologia znika.
\end{proof}
